\documentclass[12pt]{article}
\usepackage[utf8]{inputenc}
\usepackage[T1]{fontenc}
\usepackage{hyperref}
\title{Yuhup Language Resources for K-12 Curriculum Development}
\author{Compiled by Claude Sonnet 5.5 (claude-sonnet-5-5) for Andr\'as Kornai}
\date{}
\begin{document}
\maketitle

\section{Introduction}
Yuhup (ISO 639-3: yab; Glottocode yuhu1238) is a Naduhup (Nadahup, Makú) language of Brazil with about 500 speakers. This report is based on a small number of web searches. Little material exists that is specific to Yuhup, and none of it is a school resource.

\section{Primary Educational Resources}
No primers, textbooks or school materials for Yuhup were found in the search results.

\section{Linguistic and Reference Materials}
\begin{itemize}
\item Patience Epps, \textit{A Grammar of Hup}. This is a reference grammar of the closely related Hup language, not of Yuhup. Search results report that Hup and Yuhup share about 90\% of their vocabulary and are described as mutually intelligible. \url{https://shop.falter.at/detail/9783110195880/a-grammar-of-hup}
\item A dissertation version of the Hup grammar, available online. \url{https://libraetd.lib.virginia.edu/public_view/ns064640z}
\item Overview of the Nadahup languages. \url{https://en.wikipedia.org/wiki/Nadahup_languages}
\item Hunter-Gatherer language database entry, which holds a basic vocabulary list for Yuhup. \url{https://huntergatherer.la.utexas.edu/languages/language/123}
\end{itemize}
The search summary also mentioned a 2002 Paris 7 thesis by Ana Maria Ospina Bozzi on Yuhup morphology and syntax, and a 2012 study by Silva and Silva, \textit{A l\'ingua dos Yuhupdeh}. No URLs for these were seen, so they are not listed as resources here.

\section{Religious and Community Resources}
Nothing was found.

\section{Digital and Audio Resources}
Nothing Yuhup-specific was found. No dictionary was located.

\section{Curriculum Development Context}
The evidence is thin. The available material is academic: a Hup grammar and a vocabulary list. A curriculum would need new work with the community, starting from the Hup grammar and the Yuhup thesis literature. No existing school use of Yuhup was found.
\end{document}